\documentclass[10pt,letterpaper]{article}
\usepackage{cornell-notes}

\title{Annex F.04: 04-compliance-with-adf-requirements}
\author{}
\date{}
\setDocTitle{Annex F.04: 04-compliance-with-adf-requirements}
\setDocSubtitle{ISO/IEC/IEEE 42010:2022}
\setDocOwner{ISO 42010 Cornell Notes}
\setDocDate{\today}
\setCornellCollection{ISO/IEC/IEEE 42010:2022 Cornell Notes}
\setCornellUnitType{Annex}
\setCornellUnitNumber{F.04}
\setCornellUnitTitle{04-compliance-with-adf-requirements}
\hypersetup{pdftitle={Annex F.04: 04-compliance-with-adf-requirements},pdfsubject={Cornell notes},pdfauthor={}}

\begin{document}
\maketitle
\begin{CornellOverviewBox}{Overview}
{\Large\bfseries\textcolor{CNavy}{Annex F.4: Compliance with ADF requirements}}\\[3pt]
{\small\textbf{Classification:} Informative annex}\\
{\small\textbf{Learning objective:} Use the informative guidance on compliance with adf requirements to reinforce the normative and conceptual material without treating the guidance itself as mandatory.}
\end{CornellOverviewBox}
\vspace{3pt}

\begin{CornellNotesTable}
\CornellNoteRow{Purpose}{Compares several published architecture description frameworks against the framework requirements and highlights areas of full, partial, or absent coverage.}
\CornellNoteRow{Core details}{\begin{itemize}[leftmargin=*,nosep,topsep=1pt]
\item Currently none of the ADF fully comply with all the requirements itemized in 7.1.
\item Nevertheless, each of the items is addressed by several of the ADF and all these items are considered as necessary in an AD.
\item Consequently, the provisions in 7.1 provide areas for improvement for these frameworks.
\end{itemize}}
\CornellNoteRow{Comparison meaning}{The tables use “Yes” when a framework requirement is met, “Partial” when it is addressed to some degree without explicit wording, and “No” when it is not addressed.}
\CornellNoteRow{Key finding}{None of the surveyed frameworks fully covers every framework requirement; the requirement set therefore exposes concrete opportunities for framework improvement.}
\CornellNoteRow{Frameworks compared}{The comparison covers GERA, RM-ODP, Zachman, TOGAF, UAF, NAF, DoDAF, and ArchiMate.}
\CornellNoteRow{Study relationship / visual insight}{The compliance tables compare GERA, RM-ODP, Zachman, TOGAF, UAF, NAF, DoDAF, and ArchiMate across framework-identification, stakeholder/concern coverage, aspects, perspectives, formalism, viewpoints, model kinds, correspondence/legend support, and framework methods.}
\CornellNoteRow{Interpretive note}{\begin{itemize}[leftmargin=*,nosep,topsep=1pt]
\item In Table F.1 and Table F.2 “Partial” means that the requirement is fulfilled to some degree but is not expressed with explicit wording. “Yes” is used when the requirement is met. “No” states that the requirement is not addressed.
\end{itemize}}
\CornellNoteRow{Related sections}{7.1}
\CornellNoteRow{Active recall}{\begin{itemize}[leftmargin=*,nosep,topsep=1pt]
\item What is the central role of Compliance with ADF requirements?
\item How does Compliance with ADF requirements connect to adjacent architecture-description concepts?
\end{itemize}}
\end{CornellNotesTable}


\begin{CornellSummaryBox}{Summary}
Compares several published architecture description frameworks against the framework requirements and highlights areas of full, partial, or absent coverage. This material is informative guidance and does not itself create a conformance requirement.
\end{CornellSummaryBox}

\end{document}
